\chapter{Iteration 5 --- classical spectral EEG features}
\label{ch:iter5}

\section{Task definition and inputs}

\paragraph{Task.} Auditory-attention decoding: for a given trial, predict
which of four concurrently presented speakers the subject was attending
to. A trial is 30\,s long; during it the subject hears four speech
streams simultaneously (Device-1 L, Device-1 R, Device-2 L, Device-2 R
--- see Dataset Layout chapter) plus noise from Device-3 at ±135°. The
attended-speaker label comes from \code{trials.csv} and is one of
$\{1, 2, 3, 4\}$.

\paragraph{Three label-collapses evaluated.}
\begin{itemize}
  \item \textbf{Hemisphere} (chance $= 0.5$): $\{1,2\}$ (LEFT,
        az $\in\{-67.5°, -22.5°\}$) vs $\{3,4\}$ (RIGHT, az $\in\{+22.5°,
        +67.5°\}$). This is the classical binary AAD framing --- which
        side is the listener attending.
  \item \textbf{Inner vs outer} (chance $= 0.5$): $\{2,3\}$ (inner
        speakers at $\pm 22.5°$) vs $\{1,4\}$ (outer speakers at
        $\pm 67.5°$). Tests whether the decoder picks up \emph{eccentricity}
        rather than hemisphere.
  \item \textbf{4-class speaker identity} (chance $= 0.25$): predict the
        exact speaker label $1/2/3/4$.
\end{itemize}

\paragraph{Inputs used in iter-5.}
\begin{itemize}
  \item \textbf{EEG-only condition}: 32-channel EEG
        (\code{experiment\_data/Eval-K/eeg\_data.p}), bandpassed 1--40\,Hz,
        notched at 60\,Hz, auto-referenced (linked mastoids when good;
        common-average otherwise), at the native \SI{500}{\Hz}. Per trial
        we compute a \textbf{368-dimensional spectral feature vector}
        described in the next section and use it as input $X$ to a
        per-subject classifier.
  \item \textbf{Gaze-only condition} (for the fusion comparison): the
        Tobii Glasses~3 \code{gazedata.gz} stream from the Video-Recordings
        folder paired with the trial, plus the paired \code{imudata.gz}.
        The 30-s stream is summarised into a 23-dimensional vector
        (Table~\ref{tab:gaze-feats}). The experiment-data 2-D gaze
        (\code{gaze\_data.p}) is not used.
\end{itemize}

\begin{table}[ht]
\centering
\caption{Gaze / head-motion feature vector (23 features per trial).}
\label{tab:gaze-feats}
\begin{tabular}{llll}
\toprule
\# & Name & Source & Description \\
\midrule
1  & \code{gx\_mean}      & Tobii gaze2d & scene horizontal gaze, mean \\
2  & \code{gx\_std}       & Tobii gaze2d & scene horizontal gaze, std \\
3  & \code{gy\_mean}      & Tobii gaze2d & scene vertical gaze, mean \\
4  & \code{gy\_std}       & Tobii gaze2d & scene vertical gaze, std \\
5  & \code{gx\_valid}     & Tobii gaze2d & fraction of finite samples \\
6  & \code{sacc\_rate}    & derived       & saccades/s (I-VT, 30°/s thresh.) \\
7  & \code{sacc\_amp\_med} & derived      & median saccade amplitude (°) \\
8  & \code{g3d\_x\_mean}  & Tobii gaze3d & 3-D gaze x (horizontal, mm) \\
9  & \code{g3d\_y\_mean}  & Tobii gaze3d & 3-D gaze y (vertical, mm) \\
10 & \code{g3d\_z\_mean}  & Tobii gaze3d & 3-D gaze z (depth, mm) \\
11 & \code{g3d\_x\_std}   & Tobii gaze3d & 3-D gaze x, std \\
12 & \code{L\_az\_mean}   & Tobii eyeleft  & left-eye azimuth mean (°) \\
13 & \code{L\_az\_std}    & Tobii eyeleft  & left-eye azimuth std \\
14 & \code{L\_el\_mean}   & Tobii eyeleft  & left-eye elevation mean (°) \\
15 & \code{L\_pup}        & Tobii eyeleft  & left pupil diameter (mm) \\
16 & \code{L\_valid}      & Tobii eyeleft  & left-eye sample validity \\
17 & \code{R\_az\_mean}   & Tobii eyeright & right-eye azimuth mean \\
18 & \code{R\_az\_std}    & Tobii eyeright & right-eye azimuth std \\
19 & \code{R\_el\_mean}   & Tobii eyeright & right-eye elevation mean \\
20 & \code{R\_pup}        & Tobii eyeright & right pupil diameter \\
21 & \code{R\_valid}      & Tobii eyeright & right-eye sample validity \\
22 & \code{gyro\_mag}     & imudata       & mean $\|\boldsymbol{\omega}\|$ (rad/s) \\
23 & \code{acc\_std}      & imudata       & std of $\|\mathbf{a}\|-g$ (m/s$^2$) \\
\bottomrule
\end{tabular}
\end{table}

\begin{itemize}
  \setcounter{enumi}{2}
  \item \textbf{Fusion condition}: concatenation $[\text{EEG}_{368},
        \text{gaze}_{20}]$.
  \item \textbf{Scene video (\code{scenevideo.mp4})} is \emph{not} used
        in iter-5. Planned for iter-6 as motion-energy / optical-flow
        features.
\end{itemize}

\paragraph{Prediction target.} $\hat{y} = f(X)$ where $X$ is the
per-trial feature vector above and $\hat{y}$ is the attended-speaker
label under whichever of the three label-collapses we are evaluating.

\paragraph{Evaluation protocol.} For every subject independently, trials
are split with 5-fold stratified cross-validation (100 trials / subject,
so $\sim 80$ train + $\sim 20$ test per fold). The classifier is
\textbf{refit from scratch on each subject} --- no cross-subject
transfer. The reported accuracy per subject is the mean over the 5 held
-out folds; the reported \emph{pooled} accuracy is the unweighted mean
across the 16 per-subject accuracies (with bootstrap 95\,\% CI resampling
the 16 values). Chance is determined by the label collapse.

\section*{Headline}

Replacing the envelope-reconstruction decoder (CCA mel-28) with a
368-dimensional \emph{linear} classifier over classical EEG features ---
band powers, alpha lateralisation, Hjorth complexity, spectral entropy,
$1/f$ slope, alpha-band phase-locking and magnitude-squared coherence ---
takes pooled within-subject hemisphere AAD from $0.498$ to
$\mathbf{0.716}$, a $+21.8$\,pp jump. No deep learning. Subject~3 climbs
from $0.43$ (CCA) to $\mathbf{0.94}$ logreg; Subject~15 hits $\mathbf{0.92}$.

\section{Feature set}

Per trial, computed on 1--40\,Hz preprocessed, auto-referenced EEG (500\,Hz)
and aggregated over the 30\,s trial:

\begin{table}[ht]
\centering
\begin{tabular}{lrl}
\toprule
Group & \# feats & Contents \\
\midrule
Band powers (log) & 160 & Welch $\delta/\theta/\alpha/\beta/\gamma$ per 32 channels \\
Within-channel band ratios (log) & 96 & $\alpha/\theta$, $\alpha/\beta$, $\theta/\beta$ per channel \\
Alpha lateralisation & 4 & log R/L for parietal, central, occipital, temporal \\
Hjorth mobility + complexity & 64 & Per channel \\
Spectral entropy & 32 & Shannon, 1--40\,Hz \\
$1/f$ aperiodic slope & 4 & Per channel group (frontal/central/parietal/occipital) \\
Alpha connectivity & 6 & MSC \& wPLI for within-L, within-R, between-hemi \\
Global field power & 2 & mean, std \\
\midrule
\textbf{Total} & \textbf{368} & \\
\bottomrule
\end{tabular}
\end{table}

\section{Classifier sweep (all 16 subjects, per-subject 5-fold CV)}

\begin{table}[ht]
\centering
\begin{tabular}{llrl}
\toprule
Task & Classifier & Pooled acc. & 95\,\% CI \\
\midrule
Hemisphere  & L2-LogReg       & \textbf{0.716} & [0.641, 0.792] \\
Hemisphere  & LightGBM        & 0.655          & [0.589, 0.724] \\
Inner-outer & L2-LogReg       & \textbf{0.649} & [0.592, 0.712] \\
Inner-outer & LightGBM        & 0.597          & [0.536, 0.667] \\
4-class     & L2-LogReg       & \textbf{0.448} & [0.358, 0.544] \\
4-class     & LightGBM        & 0.410          & [0.332, 0.497] \\
\bottomrule
\end{tabular}
\end{table}

L2-LogReg beats LightGBM on every task --- consistent with a linearly
separable feature set and $p=368$ vs $n\approx 100$ where tree-boosting
overfits. The hemisphere pool at 0.716 is $+21.8$\,pp above the linear
TRF baseline ($0.498$, Chapter~\ref{ch:iter4}) and $+19.3$\,pp above the
best iter-3 CCA backbone pooled ($0.523$).

\section{Per-subject LogReg accuracy}

\begin{table}[ht]
\centering
\begin{tabular}{lrrr}
\toprule
Subject & Hemisphere & Inner/outer & 4-class \\
\midrule
1  & 0.480 & 0.640 & 0.280 \\
2  & \textbf{0.900} & \textbf{0.850} & 0.750 \\
3  & \textbf{0.940} & 0.600 & 0.520 \\
4  & \textbf{0.880} & \textbf{0.910} & \textbf{0.810} \\
5  & 0.700 & 0.640 & 0.460 \\
6  & 0.730 & 0.580 & 0.420 \\
7  & \textbf{0.890} & 0.730 & 0.620 \\
8  & 0.620 & 0.540 & 0.310 \\
9  & 0.600 & 0.570 & 0.240 \\
10 & 0.660 & 0.610 & 0.450 \\
11 & 0.430 & 0.510 & 0.200 \\
12 & 0.600 & 0.570 & 0.330 \\
13 & 0.770 & 0.620 & 0.400 \\
14 & 0.610 & 0.530 & 0.250 \\
15 & \textbf{0.920} & \textbf{0.889} & \textbf{0.786} \\
16 & 0.720 & 0.590 & 0.340 \\
\midrule
\textbf{Pooled} & \textbf{0.716} & \textbf{0.649} & \textbf{0.448} \\
Chance          & 0.500          & 0.500          & 0.250 \\
\bottomrule
\end{tabular}
\end{table}

Nine of 16 subjects pass $0.70$ hemisphere; five pass $0.88$. Subject~11
sits at $0.43$ --- consistent with its weak-signal profile on other
metrics (this subject also fails the linear TRF and has a below-chance
gaze-only score).

\section{Fusion with gaze features}

Early fusion = concatenate $[EEG_{368}, Gaze_{20}]$ and feed to
LightGBM ($n=300$, lr=0.05). 5-fold within-subject CV.

\begin{table}[ht]
\centering
\begin{tabular}{llll}
\toprule
Task & EEG spec. only & Gaze only & Early fusion \\
\midrule
Hemisphere  & 0.716 [0.639, 0.789] & \textbf{0.772} [0.687, 0.851] & 0.753 [0.660, 0.838] \\
Inner-outer & 0.650 [0.593, 0.714] & 0.683 [0.597, 0.766] & 0.660 [0.583, 0.735] \\
4-class     & 0.442 [0.354, 0.535] & \textbf{0.547} [0.429, 0.664] & 0.513 [0.405, 0.623] \\
\bottomrule
\end{tabular}
\end{table}

Pooled fusion does \emph{not} exceed gaze-only --- the two modalities
have substantial overlapping information because both exploit the
same underlying spatial-attention process (lateralised gaze, lateralised
alpha). \textbf{But on a per-subject basis, fusion does help when both
modalities are strong}: S7 hemisphere $0.89 \to 0.97$, S15 $0.92 \to
1.000$, S2 $0.91 \to 0.95$, S13 $0.77 \to 0.88$.

\section{What iter-5 settles}

\begin{enumerate}
  \item \textbf{Linear separation of attended hemisphere from EEG alone
        is achievable}: LogReg on 368 spectral/connectivity features
        reaches 0.716 pooled. The envelope-tracking paradigm is simply
        the wrong question to ask on this dataset.
  \item \textbf{Per-subject heterogeneity is preserved}: subjects with
        hardware-quality issues (S11, S1) remain at or below chance
        even with the richer feature set. The featurisation is not
        fixing the data; it is exposing more of the signal that was
        there all along.
  \item \textbf{Fusion gain exists per-subject but is hidden by pooled
        averaging}. On subjects where both modalities are strong, early
        fusion improves on gaze-only; on subjects where either is weak,
        it tracks the stronger modality. The pooled number is the wrong
        summary when the two modalities correlate.
  \item \textbf{LightGBM underperforms LogReg} here despite having more
        capacity --- characteristic of high-dimensional, low-noise
        linear features with only $\sim 100$ examples per subject.
        Regularised linear is the right default for this dataset size.
\end{enumerate}

\section{Iter-6 follow-ups}

\begin{enumerate}
  \item \textbf{Feature-family ablation}: drop each group one at a time
        and re-fit, to identify which subset of the 368 features
        carries the signal (expect alpha lateralisation + band powers
        in parietal/temporal channels to dominate).
  \item \textbf{Dimension reduction}: $\ell_1$-regularised logistic or
        group-sparse regression to pick an interpretable subset (goal:
        $\leq 20$ features that retain $\geq 0.70$ hemisphere).
  \item \textbf{Deep EEG model} (the planned iter-6 GPU experiment):
        EEGNet and a small transformer on 5--10\,s windows, comparing
        to the 368-feature LogReg ceiling on the same tasks.
  \item \textbf{Per-subject quality score} combining bad-channel rate,
        alpha-lat.\ direction, gaze validity, and now the iter-5
        spectral accuracy --- predict which subjects will gain from a
        deep model.
\end{enumerate}
